\documentclass[11pt, a4paper]{article}

% bfw-style.tex — gemeinsame Style-Definitionen für alle Handouts/Aufgabenhefte
% Voraussetzung: Kompilation mit xelatex (wegen fontspec)

\usepackage[ngerman]{babel}
\usepackage{geometry}
\geometry{a4paper, top=2.2cm, bottom=2.2cm, left=2.3cm, right=2.3cm}

\usepackage{fontspec}
% Fallback-Kette: Source Sans Pro -> Calibri -> Segoe UI -> Latin Modern Sans
% (Latin Modern ist immer mit MiKTeX vorhanden)
\IfFontExistsTF{Source Sans Pro}{%
  \setmainfont{Source Sans Pro}[Ligatures=TeX]%
}{\IfFontExistsTF{Calibri}{%
  \setmainfont{Calibri}[Ligatures=TeX]%
}{\IfFontExistsTF{Segoe UI}{%
  \setmainfont{Segoe UI}[Ligatures=TeX]%
}{%
  \setmainfont{Latin Modern Sans}[Ligatures=TeX]%
}}}
\IfFontExistsTF{Source Code Pro}{%
  \setmonofont{Source Code Pro}[Scale=0.92]%
}{\IfFontExistsTF{Consolas}{%
  \setmonofont{Consolas}[Scale=0.92]%
}{%
  \setmonofont{Latin Modern Mono}[Scale=0.92]%
}}

\usepackage{xcolor}
% Farbpalette: hell, freundlich, modern
\definecolor{bfwPrimary}{HTML}{0EA5A4}    % Petrol/Türkis - Primär
\definecolor{bfwPrimaryDark}{HTML}{0F766E} % dunkler Petrol für Headings
\definecolor{bfwAccent}{HTML}{F59E0B}     % Amber - Akzent
\definecolor{bfwSuccess}{HTML}{10B981}    % Grün - Erfolg / Lösung
\definecolor{bfwWarn}{HTML}{F97316}       % Orange - Achtung
\definecolor{bfwInk}{HTML}{0F172A}        % fast Schwarz
\definecolor{bfwInkSoft}{HTML}{475569}    % gedämpftes Grau
\definecolor{bfwBgSoft}{HTML}{F8FAFC}     % off-white
\definecolor{bfwBgTeal}{HTML}{ECFEFF}     % helles Türkis
\definecolor{bfwBgAmber}{HTML}{FEF3C7}    % helles Gelb
\definecolor{bfwBgRose}{HTML}{FFE4E6}     % helles Rosé für Eselsbrücken

\usepackage[colorlinks=true, linkcolor=bfwPrimaryDark, urlcolor=bfwPrimaryDark]{hyperref}
\usepackage{enumitem}
\setlist[itemize]{leftmargin=*, itemsep=2pt, topsep=2pt}
\setlist[enumerate]{leftmargin=*, itemsep=2pt, topsep=2pt}

\usepackage[most]{tcolorbox}
\usepackage{booktabs}
\usepackage{tikz}
\usetikzlibrary{decorations.pathmorphing, positioning, shapes.geometric, arrows.meta}

% Merkkästchen — wichtige Definition / Kernpunkt
\newtcolorbox{merksatz}[1][]{%
  enhanced, breakable,
  colback=bfwBgTeal,
  colframe=bfwPrimary,
  boxrule=0.6pt,
  arc=4pt,
  left=10pt, right=10pt, top=8pt, bottom=8pt,
  fonttitle=\bfseries\color{bfwPrimaryDark},
  title={\faLightbulb\ Merksatz},
  #1
}

% Eselsbrücke — Mnemoniks
\newtcolorbox{eselsbruecke}[1][]{%
  enhanced, breakable,
  colback=bfwBgRose,
  colframe=bfwAccent,
  boxrule=0.6pt,
  arc=4pt,
  left=10pt, right=10pt, top=8pt, bottom=8pt,
  fonttitle=\bfseries\color{bfwWarn},
  title={Eselsbrücke},
  #1
}

% Beispiel-Box
\newtcolorbox{beispiel}[1][]{%
  enhanced, breakable,
  colback=bfwBgSoft,
  colframe=bfwInkSoft,
  boxrule=0.4pt,
  arc=3pt,
  left=10pt, right=10pt, top=8pt, bottom=8pt,
  fonttitle=\bfseries\color{bfwInk},
  title={Beispiel},
  #1
}

% Lösungs-Box (für Aufgabenhefte)
\newtcolorbox{loesung}[1][]{%
  enhanced, breakable,
  colback=bfwBgAmber!40,
  colframe=bfwSuccess,
  boxrule=0.6pt,
  arc=3pt,
  left=10pt, right=10pt, top=8pt, bottom=8pt,
  fonttitle=\bfseries\color{bfwSuccess!70!black},
  title={Lösung},
  #1
}

% Glossar-Eintrag
\newcommand{\glossareintrag}[2]{%
  \par\noindent\textbf{\color{bfwPrimaryDark}#1} \enspace #2\par\smallskip
}

% Section-Stil — etwas farbiger als Standard
\usepackage{titlesec}
\titleformat{\section}
  {\Large\bfseries\color{bfwPrimaryDark}}
  {\thesection}{0.6em}{}
\titleformat{\subsection}
  {\large\bfseries\color{bfwPrimary}}
  {\thesubsection}{0.5em}{}
\titleformat{\subsubsection}
  {\normalsize\bfseries\color{bfwInk}}
  {\thesubsubsection}{0.4em}{}

% TikZ-Stil "skizzenhaft" — etwas weicher als Rough.js, aber stabil
\tikzset{
  roughbox/.style = {
    draw=bfwPrimaryDark, thick, fill=bfwBgTeal,
    rounded corners=4pt, inner sep=6pt
  },
  roughaccent/.style = {
    draw=bfwAccent, thick, fill=bfwBgAmber,
    rounded corners=4pt, inner sep=6pt
  },
  roughline/.style = {
    draw=bfwInkSoft, thick, -{Stealth[length=2.5mm]}
  }
}

% Bullet-Symbole (bewusst kein FontAwesome — vermeidet Font-Installations-Probleme)
\newcommand{\faLightbulb}{$\bullet$}
\newcommand{\faBookmark}{$\bullet$}

% Header/Footer
\usepackage{fancyhdr}
\pagestyle{fancy}
\fancyhf{}
\renewcommand{\headrulewidth}{0pt}
\fancyfoot[L]{\small\color{bfwInkSoft}\thepage}
\fancyfoot[R]{\small\color{bfwInkSoft} BFW M-064 Beschaffung im Büromanagement}


\title{\vspace*{-1.5cm}\Huge\bfseries\color{bfwPrimaryDark}Aufgabenheft Tag~7\\[0.4em]
  \large\color{bfwInkSoft}\textnormal{Sitzungen \& Besprechungen vorbereiten und durchführen --- Übungen im IHK-Klausurformat}}
\author{\textsf{Büroprozesse gestalten und Arbeitsvorgänge organisieren \textperiodcentered{} M-067}\\
\textsf{\small\copyright\ Can Siebert\,\textperiodcentered\,2026}}
\date{Selbstlernmaterial \textperiodcentered{} 07.07.2026}

\begin{document}
\maketitle
\thispagestyle{empty}

\begin{merksatz}[title={\bfseries So nutzen Sie dieses Aufgabenheft}]
Bearbeiten Sie alle Aufgaben zunächst \emph{ohne} in die Lösungen zu schauen. Schreiben
Sie Ihre Antworten in vollständigen Sätzen --- die IHK-Prüfung verlangt formulierte
Begründungen, nicht bloße Stichworte. Beachten Sie die \emph{Operatoren} (nennen,
erläutern, beurteilen, begründen) --- sie geben den geforderten Antwort-Tiefegrad vor.
Erst danach öffnen Sie den Lösungsteil ab Seite~\pageref{loesungen} und vergleichen.
\end{merksatz}

\tableofcontents
\vspace{1em}
\hrule

\section{Aufgaben}

\subsection*{Aufgabe~1 --- Warum Vorbereitung entscheidet (10 Punkte)}

\noindent\textbf{\color{bfwAccent}YouTube-Vorbereitung:}\enspace \href{https://www.youtube.com/watch?v=DHG3tcWrwmo}{Besprechung vorbereiten --- Grundlagen} (\url{https://www.youtube.com/watch?v=DHG3tcWrwmo})

\medskip
Sie sind im Sekretariat der \emph{Duisdorfer BüroKonzept KG} tätig. Ihre Vorgesetzte klagt,
dass die letzte Besprechung „nichts gebracht“ habe, und bittet Sie um eine Einschätzung.

\begin{enumerate}[label=(\alph*)]
  \item \textbf{Erläutern} Sie, warum eine gute Vorbereitung für den Erfolg einer Sitzung unverzichtbar ist. \hfill (3~P.)
  \item \textbf{Nennen} Sie drei Folgen einer mangelhaft vorbereiteten Sitzung. \hfill (3~P.)
  \item \textbf{Erläutern} Sie, wovon der Zeitpunkt des Vorbereitungsbeginns abhängt. \hfill (2~P.)
  \item \textbf{Beurteilen} Sie die Aussage: \emph{„Für eine kleine Teamsitzung braucht man keine Vorbereitung.“} \hfill (2~P.)
\end{enumerate}

\subsection*{Aufgabe~2 --- Vorbereitungsaspekte klären (12 Punkte)}

\noindent\textbf{\color{bfwAccent}YouTube-Vorbereitung:}\enspace \href{https://www.youtube.com/watch?v=YROwKwdHYdk}{Sitzung planen --- Checkliste} (\url{https://www.youtube.com/watch?v=YROwKwdHYdk})

\medskip
Vor dem Verfassen der Einladung sind mehrere Vorbereitungsaspekte zu klären.

\begin{enumerate}[label=(\alph*)]
  \item \textbf{Nennen} Sie sechs Vorbereitungsaspekte einer Sitzung. \hfill (3~P.)
  \item \textbf{Ordnen} Sie drei davon jeweils eine passende Leitfrage \textbf{zu}. \hfill (3~P.)
  \item \textbf{Erläutern} Sie, worauf bei der Wahl von Datum und Uhrzeit zu achten ist. \hfill (3~P.)
  \item \textbf{Beurteilen} Sie, warum Raum, Technik und Verpflegung frühzeitig geklärt werden sollten. \hfill (3~P.)
\end{enumerate}

\subsection*{Aufgabe~3 --- Tagesordnung formulieren (12 Punkte)}

\noindent\textbf{\color{bfwAccent}YouTube-Vorbereitung:}\enspace \href{https://www.youtube.com/watch?v=_KZaG_BxFwo}{Tagesordnung \& TOP erstellen} (\url{https://www.youtube.com/watch?v=_KZaG_BxFwo})

\medskip
In einer Einladung ist ein Tagesordnungspunkt nur mit dem Wort \emph{„Sortiment“} bezeichnet.

\begin{enumerate}[label=(\alph*)]
  \item \textbf{Erläutern} Sie, welches Problem diese ungenaue Formulierung verursacht. \hfill (3~P.)
  \item \textbf{Formulieren} Sie den TOP so präzise um, dass sich die Teilnehmer gezielt vorbereiten können. \hfill (3~P.)
  \item \textbf{Nennen} Sie vier mögliche Ziele eines Tagesordnungspunkts. \hfill (3~P.)
  \item \textbf{Erläutern} Sie, warum eine gewisse Wiederholung im Sitzungsablauf sinnvoll ist. \hfill (3~P.)
\end{enumerate}

\subsection*{Aufgabe~4 --- Einladung erstellen (15 Punkte)}

\noindent\textbf{\color{bfwAccent}YouTube-Vorbereitung:}\enspace \href{https://www.youtube.com/watch?v=B7xyYof_S-k}{Einladung zur Besprechung schreiben} (\url{https://www.youtube.com/watch?v=B7xyYof_S-k})

\medskip
Die Geschäftsführung lädt zu einer Abteilungsleitersitzung am 30.09. um 14:00~Uhr in Raum~236
ein. Zeitbedarf ca.\ 90~Minuten, das Protokoll führt Samira Erden.

\begin{enumerate}[label=(\alph*)]
  \item \textbf{Nennen} Sie sechs Angaben, die zu einer detaillierten Einladung gehören. \hfill (6~P.)
  \item \textbf{Erstellen} Sie für diese Sitzung eine vollständige Einladung (Interne Mitteilung) mit mindestens vier Tagesordnungspunkten. \hfill (6~P.)
  \item \textbf{Begründen} Sie, warum die Einladung erst nach Abschluss der übrigen Vorbereitungen gestaltet wird. \hfill (3~P.)
\end{enumerate}

\subsection*{Aufgabe~5 --- Durchführung und Teilnehmertypen (12 Punkte)}

\noindent\textbf{\color{bfwAccent}YouTube-Vorbereitung:}\enspace \href{https://www.youtube.com/watch?v=D8QTrMrTEyM}{Besprechung moderieren --- Regeln} (\url{https://www.youtube.com/watch?v=D8QTrMrTEyM})

\medskip
Trotz guter Vorbereitung verläuft eine Sitzung zäh: Eine Person redet ständig, eine andere
schaut nur auf die Uhr.

\begin{enumerate}[label=(\alph*)]
  \item \textbf{Nennen} Sie vier Grundsätze für die Durchführung einer Besprechung. \hfill (4~P.)
  \item \textbf{Beschreiben} Sie die drei im Lehrbuch genannten schwierigen Teilnehmertypen. \hfill (3~P.)
  \item \textbf{Erläutern} Sie für zwei dieser Typen je eine passende Reaktion der Moderation. \hfill (3~P.)
  \item \textbf{Beurteilen} Sie die Aussage: \emph{„Für das Gelingen ist allein die Leitung verantwortlich.“} \hfill (2~P.)
\end{enumerate}

\subsection*{Aufgabe~6 --- Elektronische und hybride Besprechungen (10 Punkte)}

\noindent\textbf{\color{bfwAccent}YouTube-Vorbereitung:}\enspace \href{https://www.youtube.com/watch?v=0FnoyDzE4z8}{Videokonferenz --- Etikette \& Tipps} (\url{https://www.youtube.com/watch?v=0FnoyDzE4z8})

\medskip
\begin{enumerate}[label=(\alph*)]
  \item \textbf{Erläutern} Sie, was eine Videokonferenz ermöglicht und welche Voraussetzungen dafür nötig sind. \hfill (3~P.)
  \item \textbf{Erklären} Sie den Unterschied zwischen einer Telefonkonferenz und einer Dreierkonferenz. \hfill (2~P.)
  \item \textbf{Nennen} Sie zwei Vorteile elektronischer Besprechungen. \hfill (2~P.)
  \item \textbf{Nennen} Sie drei Etikette-Regeln für eine Videokonferenz und \textbf{erläutern} Sie, worauf bei Hybridmeetings besonders zu achten ist. \hfill (3~P.)
\end{enumerate}

\vspace{1em}
\noindent\textbf{Punkte gesamt: 71}

\newpage

\section{Lösungen}\label{loesungen}

\subsection*{Lösung Aufgabe~1}
\begin{loesung}
\textbf{(a)} Ob eine Sitzung zufriedenstellend und effektiv verläuft, entscheidet sich
meist schon bei der Vorbereitung. Ohne Planung droht, dass wichtige Inhalte nicht
thematisiert werden, die Sitzung verspätet beginnt, mehr Zeit benötigt und Ergebnisse
nicht festgehalten werden. Gute Vorbereitung spart somit Zeit und sichert Ergebnisse.

\textbf{(b)} Drei Folgen (Beispiele): wichtige Themen werden nicht behandelt; verspäteter
Beginn und höherer Zeitbedarf; Beschlüsse werden nicht genau festgehalten, nicht
nachverfolgt und schließlich vergessen.

\textbf{(c)} Der Zeitpunkt hängt von der \emph{Art der Sitzung} und den konkreten
\emph{Bedingungen vor Ort} ab: Eine wöchentliche Teamsitzung lässt sich 2--3 Tage vorher
planen, eine offizielle Sitzung (z.\,B.\ Gesellschafterversammlung) erfordert wegen
gesetzlicher Fristen einen früheren Beginn.

\textbf{(d)} Die Aussage ist falsch. Eine gute Planung ist auch für Teambesprechungen mit
kleiner Teilnehmerzahl wichtig --- sonst werden Themen vergessen, es geht Zeit verloren und
Ergebnisse bleiben unverbindlich.
\end{loesung}

\subsection*{Lösung Aufgabe~2}
\begin{loesung}
\textbf{(a)} Vorbereitungsaspekte: \textbf{Datum und Uhrzeit, Raum, Teilnehmer, Zeitbedarf,
Tagesordnungspunkte (TOP), Unterlagen, Protokoll}.

\textbf{(b)} Leitfragen (Beispiele):
\begin{itemize}[itemsep=0pt]
  \item \emph{Datum/Uhrzeit}: Gibt es andere Sitzungen zum Termin? Sind Teilnehmer verhindert?
  \item \emph{Raum}: Ist der Raum geeignet, verfügbar und reserviert? Welche Technik wird benötigt?
  \item \emph{Protokoll}: Wer schreibt das Protokoll? Welche Protokollart ist gewünscht?
\end{itemize}

\textbf{(c)} Bei Datum und Uhrzeit ist zu prüfen, ob es konkurrierende Sitzungen zum
geplanten Termin gibt und ob wichtige Teilnehmer eventuell verhindert sind --- nur so sind
alle Benötigten anwesend.

\textbf{(d)} Raum, Technik und Verpflegung müssen frühzeitig geklärt werden, weil ein
passender Raum reserviert, technische Hilfsmittel (Beamer, PC) organisiert und Getränke
bereitgestellt werden müssen. Fehlt dies, kann die Sitzung nicht wie geplant stattfinden
und Zeit geht verloren.
\end{loesung}

\subsection*{Lösung Aufgabe~3}
\begin{loesung}
\textbf{(a)} Mit „Sortiment“ bleibt unklar, was besprochen werden soll (Erweiterung,
Herausnahme von Produkten oder Neuzusammenstellung). Die Teilnehmer können sich gedanklich
nicht gezielt vorbereiten und gehen ohne durchdachte Vorschläge oder Ergebnisse in die
Sitzung.

\textbf{(b)} Präzise Formulierung (Beispiel): \emph{„Aussprache und Entscheidung über die
Sortimentserweiterung der Warengruppe Handelswaren Bürotechnik (Anlagen 1 und 2)“}.

\textbf{(c)} Mögliche Ziele eines TOP: Information, Ideensammlung, Erfahrungsaustausch,
Aussprache, Entscheidungsfindung, Beschlussfassung (vier genügen).

\textbf{(d)} Eine gewisse Wiederholung schafft Routine und Orientierung: Immer
wiederkehrende Punkte (z.\,B.\ „Genehmigung des Protokolls“, „Verschiedenes“) erhalten einen
festen Platz im Ablauf, was die Vorbereitung und Durchführung erleichtert.
\end{loesung}

\subsection*{Lösung Aufgabe~4}
\begin{loesung}
\textbf{(a)} Angaben einer detaillierten Einladung: Art der Sitzung; Datum, Uhrzeit und Ort;
voraussichtlicher Zeitbedarf; Tagesordnungspunkte mit Ziel; Berichterstatter/Verantwortliche
je TOP; Hinweis auf Anlagen; Hinweis, ob Anmeldung erforderlich ist; Name des
Protokollführers.

\textbf{(b)} Beispiel-Einladung:
\begin{itemize}[itemsep=0pt]
  \item \emph{Art:} Einladung zur Abteilungsleitersitzung
  \item \emph{Datum/Ort:} 30.09., 14:00~Uhr, Raum~236 \textbar{} \emph{Zeitbedarf:} ca.\ 90~Minuten \textbar{} \emph{Protokoll:} Samira Erden
  \item \emph{Tagesordnung:} 1.\ Genehmigung des Protokolls der letzten Sitzung; 2.\ Kurzberichte aus den Abteilungen (alle Abteilungsleiter/-innen); 3.\ Aussprache und Entscheidung über die Sortimentserweiterung Handelswaren Bürotechnik (Anlagen 1 und 2, Berichterstattung: K.\ Niester); 4.\ Vorschläge für den Betriebsausflug; 5.\ Ausblick auf Oktober; 6.\ Verschiedenes.
\end{itemize}

\textbf{(c)} Die Einladung wird erst nach Abschluss der Vorbereitungen gestaltet, weil erst
dann die Inhalte präzise formuliert werden können. Je präziser die Einladung, desto besser
können sich die Teilnehmer vorbereiten und Mitverantwortung für eine effektive Durchführung
übernehmen.
\end{loesung}

\subsection*{Lösung Aufgabe~5}
\begin{loesung}
\textbf{(a)} Grundsätze (vier von): pünktlicher Beginn und Einhalten des Zeitrahmens; die
Tagesordnung als Leitfaden, zielorientiert und ohne Binnengespräche arbeiten; strukturierter
Ablauf je TOP (fragen -- zuhören -- klären -- festhalten); offene, respektvolle Kommunikation,
dem Redner nicht ins Wort fallen, Konsens suchen; jeder bemüht sich um kurze Beiträge.

\textbf{(b)} Teilnehmertypen: der \emph{Dauerredner}, der zu allem etwas zu sagen hat; der
\emph{Desinteressierte}, der immer wieder auf die Uhr schaut; der \emph{Abteilungswitzbold},
der auf Kosten anderer Scherze macht und Beiträge ins Lächerliche zieht.

\textbf{(c)} Reaktionen (zwei von): Dauerredner freundlich begrenzen und auf Zeit/Thema
verweisen; Desinteressierten gezielt einbinden und nach seiner Meinung fragen; Witzbold
sachlich zum Thema zurückführen und Beiträge wertschätzen.

\textbf{(d)} Die Aussage ist falsch. Für das gute Gelingen ist nicht nur die Leitung,
sondern sind \emph{auch die Teilnehmer} verantwortlich --- etwa durch Pünktlichkeit, kurze
Beiträge und respektvolle Kommunikation.
\end{loesung}

\subsection*{Lösung Aufgabe~6}
\begin{loesung}
\textbf{(a)} Eine Videokonferenz ermöglicht, dass sich Gesprächspartner an weit
auseinanderliegenden Orten austauschen und gemeinsam an Dateien arbeiten bzw.\ Dateien
austauschen. Voraussetzungen sind ein spezielles Videokonferenzsystem und ein schneller
Internetanschluss.

\textbf{(b)} Bei einer Telefonkonferenz wird ein Telefonat mit \emph{mehr als drei}
Teilnehmern geführt; sind es \emph{genau drei}, spricht man von einer Dreierkonferenz.

\textbf{(c)} Vorteile (zwei von): Einsparung von Kosten (Flug-, Hotelkosten, Verpflegung)
und Zeit; ortsunabhängige Zusammenarbeit; gemeinsames Arbeiten an Dateien.

\textbf{(d)} Etikette-Regeln (drei von): Technik vorher testen und pünktlich erscheinen;
Mikrofon stummschalten, wenn man nicht spricht; Kamera einschalten; ruhiger, aufgeräumter
Hintergrund; nur eine Person spricht zur Zeit. Bei \emph{Hybridmeetings} ist besonders darauf
zu achten, die online zugeschalteten Teilnehmer aktiv und gleichberechtigt einzubinden, damit
sie nicht benachteiligt werden.
\end{loesung}

\end{document}
